% !TEX root = ../Main.tex
\chapter{自由泳}

\section{泳式选择}

\state{何谓"自由"}{
    \begin{itemize}
        \item 在\textbf{自由泳比赛}中，可采用\textbf{任何泳式}；
        \item 但在\textbf{个人混合泳}及\textbf{混合泳接力}比赛中，"自由泳" 指的是\textbf{除蝶泳、仰泳、蛙泳以外的泳式}（即只能用"第 4 种"姿势）。
    \end{itemize}
}

\section{比赛规则}

\state{触壁与姿势}{
    \begin{itemize}
        \item 每次\textbf{转身}和\textbf{到达终点}时，运动员\textbf{身体的某一部分必须触及池壁}；
        \item 在\textbf{整个游程中}，运动员\textbf{身体的某一部分必须露出水面}。
    \end{itemize}
}

\state{15 米规则}{
    \begin{itemize}
        \item 在\textbf{出发}和\textbf{每次转身}后，允许运动员身体\textbf{完全没入水中}；
        \item 但在\textbf{15 米前（含 15 米）}，运动员的\textbf{头部必须露出水面}。
    \end{itemize}
}

\rmkb{
    \textbf{易考点}：
    \begin{itemize}
        \item 自由泳\textbf{可采用任何泳式}——但\textbf{混合泳中的"自由泳"}不能用蝶 / 仰 / 蛙；
        \item 终点\textbf{一只手}触壁即可；
        \item \textbf{15 米规则}：出发 / 转身后最多潜游 15 米，\textbf{15 米前（含）头必须露出}。
    \end{itemize}
}
